\textbf{GCD Properties:}
In particular, recalling that gcd is a positive integer valued function we obtain that $\gcd(a, b \cdot c) = 1$ if and only if $\gcd(a, b) = 1$ and $\gcd(a, c) = 1$.

15. If none of $a_1, a_2, \ldots, a_r$ is zero, then $\gcd(a_1, a_2, \ldots, a_r) = \gcd(\gcd(a_1, a_2, \ldots, a_{r-1}), a_r)$

16. $\gcd(a,b) \cdot \text{lcm}(a,b) = |a \cdot b|$, so for any two numbers LCM is always multiple of GCD.

17. In a Cartesian coordinate system, $\gcd(a, b)$ can be interpreted as the number of segments between points with integral coordinates on the straight line segment joining the points $(0, 0)$ and $(a, b)$.

18. Two numbers are called relatively prime, or coprime, if their greatest common divisor equals 1. For example, 9 and 28 are relatively prime.

19. Any two consecutive integers will be Co-Prime. Ex: 25,26 is co-prime. 5,6 is co-prime and so on. This will work because, let x divides 25 so remainder is something. If we divided 26 by x then remainder will be 1. This is the logic.

20. Two or more odd consecutive positive integer GCD is also 1.

\begin{lstlisting}
int gcdFast(int a, int b)
{
  int result;
  __asm__ __volatile__(
    "movl %1, %%eax;"
    "movl %2, %%ebx;"
    "CONTD: cmpl $0, %%ebx;"
    "je DONE;"
    "xorl %%edx, %%edx;"
    "idivl %%ebx;"
    "movl %%ebx, %%eax;"
    "movl %%edx, %%ebx;"
    "jmp CONTD;"
    "DONE: movl %%eax, %0;"
    : "=g"(result)
    : "g"(a), "g"(b));
  return result;
}
\end{lstlisting}

\textbf{Extended Euclid and CRT}
\begin{lstlisting}
using T = __int128_t;
T egcd(T a, T b, T &x, T &y) {
  T xx = y = 0, yy = x = 1;
  while (b) {
    T q = a / b;
    T t = b; b = a % b; a = t;
    t = xx; xx = x - q * xx; x = t;
    t = yy; yy = y - q * yy; y = t;
  }
  return a;
}
// finds x such that x % m1 = a1, x % m2 = a2.
// m1 and m2 may not be coprime
// here, x is unique modulo m = lcm(m1, m2).
// returns (x, m). on failure, m = -1.
pair<T, T> crt(T a1, T m1, T a2, T m2) {
  T p, q;
  T g = egcd(m1, m2, p, q);
  if (a1 % g != a2 % g) return {0, -1};
  T m = m1 / g * m2;
  p = (p % m + m) % m;
  q = (q % m + m) % m;
  T x = p * a2 % m * (m1 / g) % m + q * a1 % m *
  (m2 / g) % m;
  return {x % m, m};
}
\end{lstlisting}

\textbf{SPF, Mobius, Primes in O(N)}
\begin{lstlisting}
const int N = 1e6 + 2;
vector<int> spf(N), mob(N), primes;
void build() {
  mob[1] = 1;
  for (int i = 2; i < N; i++) {
    if (spf[i] == 0) {
      spf[i] = i, mob[i] = -1,
      primes.push_back(i);
    }
    for (int j = 0; i * primes[j] < N; j++) {
      spf[i * primes[j]] = primes[j];
      if (i % primes[j])
        mob[i * primes[j]] = mob[i] *
        mob[primes[j]];
      else mob[i * primes[j]] = 0;
      if (primes[j] == spf[i]) break;
    }
  }
}
\end{lstlisting}
